% Pursuing Stacks (cotes 134-2 to 134-8): archive cote 134-7, site page 35 (PDF page 36), typescript page 477.
% typescript. Draft compared: ps6.tex:1668-1740.
% AI-transcribed by Claude Opus 5.5 (claude-opus-5-5, adaptive thinking, effort high), Message Batches API, batch ps3 (msgbatch_01XdFmuVVBAcgc7EFxdtCqsj), 2026-11-12; not yet checked by eye.
\scanpage{134-7/35}
\mspage{477}
fact,  ring extension for homotopy types (an operation of greater interest
than ring restriction surely) is \uline{not} expressed,  in general,  by just
performing the trivial ring extension operation
\[ X_{*} \mapsto X_{*} \otimes_{k} k' \]
on models in $\hand{\unsure{\pi|\mathcal{M}}}_{n}(k)$,  unless we assume k' to be flat over k say  -  but
even in this case it is by no means clear a priori that the operation
above transforms weak equivalences into weak equivalences.  This is very
clearly \deleted{suggested}\inserted{shown} by the linear analogon,  the categories $\hand{\mathcal{M}}_{n}(k)$  being
replaced by \inserted{the} categories of chain complexes in  $Ab_{k}$ say,  or by $Comp^{-}(Ab_{k})$
or the like.  \xout{In} In order to correctly describe ground ring extension on
homotopy types,  we'll have first to take a suitable``resolution'' of $X_{*}$ ,
namely replace $X_{*}$  by some  $K_{*}$  say,  endowed with a weak equivalence
\[ K_{*} \to X_{*} \quad , \]
and $K_{*}$ satisfying some extra assumptions.  Maybe flatness of the compo-%
nent would be enough here.  For other purposes,  we may have to use
resolutions which are even smooth (componentwise),  or which satisfy a
suitable \uline{Kan condition}  (of a type outlined on page 463),  or both.
Our expectation is that,  when we restrict to the subcategory
\[ \deleted{\unsure{SK_{n}(k)}} \ \hand{sK\mathcal{M}_{n}(k)} \]
of the model category $\hand{\mathcal{M}}_{n}(k)$ made up with smooth Kan complexes,  \uline{that}
\uline{the category} $Hot_{n}(k)$  \uline{may be described simply in  terms of such}  sK-%
\uline{complexes} \uline{``up to h\hand{$\smile$}omotopy''},  as usual.  If this is so,  the ground ring
extension functor \deleted{(S)}  follows trivially from a corresponding functor on
the \xout{\unsure{Kss}} sK-model categories
\[ \text{(T)} \qquad sK\hand{\mathcal{M}}_{n}(k) \longrightarrow sK\hand{\mathcal{M}}_{n}(k') \quad , \quad \hand{X_{*} \longmapsto X_{*} \otimes_{k} k'} \hand{,} \]
hence
\[ \text{(T')} \qquad Hot_{n}(k) \longrightarrow Hot_{n}(k') \quad . \]

From the \inserted[typed]{sections} functor  (R')  we get homotopy invariants  $\hand{\pi}_{i}$ for an object
in  $Hot_{n}(k)$,  but the relevant k-module structure on these is not appea-%
rant \unsure{on|in} this definition.  We have a better hold,  via linear algebra over
k,  upon  homology invariants of $X_{*}$  ,  which are k-linear objects,  and
are definitely distinct (unless $k=\underline{Z}$ )  from the homology invariants of
$X_{*}(k)$ ,  which for gene ral k are definitely of little interest it seems.
The definition of homology goes via the pointed linearization functor \hand{(M)}\hand{:}
\[ \text{(U)} \qquad \hand{\Big\{} \begin{array}{l} LH.(X_{*}) \overset{def}{=} \xout{\unsure{KK}} \ L_{pt}\hand{(X_{*})} \inserted{, viewed as an object in $D.(Ab_{k})$,} \\ \hand{H_{i}(X_{*}) \overset{def}{=} H_{i}(LH.(X_{*})) = \pi_{i}(L_{pt}(X_{*}))} \ \hand{in $Ab_{k}$,} \end{array} \]
\margin{35}
